% TOP_PROBLEM: {"id": 1248, "title": "Infinitely Many Perfect Numbers", "category_id": 1, "category": {"id": 1, "name": "number_theory", "display_name": "Number Theory", "description": "Properties of integers, prime numbers, Diophantine equations.", "slug": "number-theory", "order_index": 1, "created_at": "2026-07-31T15:26:25.670Z"}, "status": "open"}
% FIRST_POSTED: null
% ENTRY_KIND: statement_only
% Imported from: batch05.tex; UnsolvedMath ID 1248
\documentclass[11pt]{article}
\usepackage[margin=25mm]{geometry}
\usepackage{fontspec}
\setmainfont{Latin Modern Roman}
\usepackage{amsmath,amssymb,mathtools}
\usepackage{unicode-math}
\setmathfont{Latin Modern Math}
\usepackage{xurl,hyperref}
\hypersetup{colorlinks=true,urlcolor=blue}
\setcounter{secnumdepth}{0}
\setlength{\parindent}{0pt}
\setlength{\parskip}{5pt}
\setlength{\emergencystretch}{3em}
\newcommand{\sref}[2]{\hyperlink{source:#1:#2}{[S#2]}}
\newcommand{\cataloguescope}[1]{\paragraph{Recorded target.}#1}
\begin{document}
% UnsolvedMath ID 1248; source catalogue code NT-041
\setcounter{section}{1247}
\section{Infinitely Many Perfect Numbers}\label{1248}
{\small\sffamily UnsolvedMath ID 1248 · Number Theory}
\subsection{Definitions and mathematical statement}

\paragraph{Shared notation.}$\mathbb N=\{1,2,\ldots\}$, and $\mathbb Z,\mathbb Q,\mathbb R,\mathbb C$ denote the integers, rationals, reals and complex numbers. Any other conventions are specified locally.


For $n\in\mathbb N$, let $\sigma(n)=\sum_{d\mid n,\ d>0}d$ be the sum of its positive divisors. A perfect number satisfies $\sigma(n)=2n$, equivalently $n$ is the sum of its positive proper divisors. Both even and odd positive integers are included.
\paragraph{Statement recorded in the supplied source.}
Is the set $\{n\in\mathbb N:\sigma(n)=2n\}$ infinite? Equivalently, for every $B\in\mathbb N$, must there exist $n>B$ with $\sigma(n)=2n$? \sref{1248}{2}
The Euclid--Euler description identifies the even perfect numbers with $2^{p-1}(2^p-1)$ for prime Mersenne numbers $2^p-1$. Infinitely many Mersenne primes would imply the target, but the recorded question does not exclude possible odd perfect numbers.
\subsection{Short English statement}
Are there infinitely many positive integers equal to the sum of their proper positive divisors, including both parity classes?
\subsection{Sources}
\begin{description}
\item[\hypertarget{source:1248:1}{S1}] ULAM.AI and the UnsolvedMath Contributors, UnsolvedMath: A Curated Collection of Open Mathematics Problems, record 1248 (NT-041). CC BY 4.0. \url{https://huggingface.co/datasets/ulamai/UnsolvedMath}.
\item[\hypertarget{source:1248:2}{S2}] John Cowles and Ruben Gamboa, \emph{Perfect Numbers in ACL2}, EPTCS 192 (2015), pp.\ 53--59, Sections 1 and 2; doi:10.4204/EPTCS.192.5. \url{https://arxiv.org/pdf/1509.06081}.
\end{description}
\label{end:1248}
\end{document}
